% ============================================================
% Preface
% Designing Agentic AI Solutions:
% Ten Business Cases with RAG, Agents, and MCP --
% from Whiteboard to Production
% Chris Kuo
% ============================================================
\chapter*{Preface}
\addcontentsline{toc}{chapter}{Preface}
\markboth{Preface}{Preface}

\section*{The question students keep asking}

Every term, my students arrive knowing more vocabulary than the term before.
They can explain retrieval-augmented generation. They have built a chatbot
that calls a tool. Many of them have heard of the Model Context Protocol and
can tell you that it is ``USB-C for AI.'' Yet when I hand them a real business
problem, such as a claims department drowning in paperwork or a retailer whose
customers cannot find answers in a thousand-page policy manual, and ask them to
\emph{design} a solution, the room goes quiet.

The quiet is not a lack of intelligence or effort. It is a gap in what we
teach. Tutorials show how to use each component. Very few show how to decide
which components a problem needs, how they should talk to each other, where
the data should live, which steps a human must approve, and what the whole
thing will cost to run on a Tuesday afternoon when traffic doubles. Those
decisions are design decisions, and design is the skill that separates a
working demo from a working business.

This book is my attempt to close that gap.

\section*{What this book is about}

This is a book about \emph{designing} agentic AI solutions. Its three building
blocks are:

\begin{itemize}
  \item \textbf{Retrieval-augmented generation (RAG)}, which grounds a model's
        answers in an organization's own documents and data;
  \item \textbf{Agents}, which plan, call tools, and carry a task across many
        steps toward a goal; and
  \item \textbf{The Model Context Protocol (MCP)}, which gives agents a standard,
        governable way to reach the systems a business already runs: databases,
        ticketing systems, file stores, and internal APIs.
\end{itemize}

Each of these is useful on its own. The interesting problems, and the hard
ones, appear when they are combined. When should a question be answered by
retrieval and when by a tool call? Should one agent do everything, or should
several specialists hand work to one another? Which capabilities belong behind
an MCP server, and who is allowed to call them? This book treats those questions
as the main subject, not as footnotes to a code listing.

\section*{Ten business cases, from whiteboard to production}

The book is built around ten business cases. I chose cases rather than topics
because design cannot be learned in the abstract. A design is only good or bad
relative to a problem: its users, its data, its risks, and its budget.

Every case follows the same arc, so that by the third or fourth case the arc
itself becomes a habit of mind:

\begin{enumerate}
  \item \textbf{The business problem.} Who is the customer, what hurts, and how
        will we know the pain went away? We define success in business terms,
        such as hours saved, errors avoided, or revenue protected, before we write
        any prompts.
  \item \textbf{The whiteboard design.} We draw the solution as a flowchart or
        graph: the agents, the retrieval sources, the MCP servers, the human
        checkpoints, and the flows of data and control between them. The
        diagram is the primary artifact of each chapter. If you cannot draw it,
        you do not yet understand it.
  \item \textbf{Trade-offs.} We weigh the pros and cons of the design against at
        least one alternative. Every choice has a price in latency, cost,
        accuracy, security, or maintainability, and a designer's job is to know
        which price is being paid and why it is worth paying.
  \item \textbf{Critique and redesign.} Several cases begin from a publicly
        available reference design, including sample agents published by
        Google Cloud. We read these samples the way an architect reads someone
        else's blueprint: respectfully, but critically. Where do they fall
        short for a real business? Then we redesign. The remaining cases are
        original designs built from a blank whiteboard.
  \item \textbf{Build.} We implement the design, using Google's Agent
        Development Kit (ADK) and Google Cloud as our working toolchain.
  \item \textbf{Production.} We take the solution the last mile: evaluation,
        observability, security and access control, cost management, and
        deployment. A design that cannot survive production was never a design
        at all.
\end{enumerate}

% TODO: insert a short table or paragraph listing the ten cases by industry
% once the case list is final.

\section*{Why the business side matters}

Students who study AI engineering often treat the business context as
background color. In practice it is the opposite: the business context is what
determines the architecture. A bank and a coffee chain may both want a
``customer service agent,'' but the bank's regulators, audit trail, and cost of
a wrong answer lead to a very different design than the coffee chain's. For this
reason every case in this book opens with the business and returns to it at the
end. My hope is that readers will finish the book able to sit across the table
from a manager, ask the right questions, and propose a solution that makes sense
on a balance sheet as well as in a code review.

\section*{A note on tools}

This book uses Google's Agent Development Kit and Google Cloud to carry each case
from whiteboard to production. They are capable, well documented, and they let
us show a complete path to deployment in one consistent environment.

But this is not a book about one vendor. Frameworks change quickly; the design
principles in these chapters do not. An agent that should hand off to a human
before it issues a refund should do so whether it runs on ADK, another
framework, or something not yet invented. Wherever possible I separate the
\emph{design decision} from the \emph{implementation detail}, so that readers
who work on another platform can carry the designs with them. MCP, as an open
protocol, helps here: a well-designed MCP server is useful to any agent that
speaks the protocol.

\section*{Who this book is for}

I wrote this book for three groups of readers.

\begin{itemize}
  \item \textbf{Students} in courses on AI engineering, data science, information
        systems, or business analytics who already know the basics of large
        language models and want to learn how to design complete solutions.
  \item \textbf{Practitioners}, including engineers, data scientists, solution
        architects, and technical product managers, who are being asked to
        ``add agents'' to their organization and want a disciplined way to
        decide what to build.
  \item \textbf{Instructors} who want a case-based structure for a course. Each
        case is designed to fill about one week of class: one session on the
        business problem and the whiteboard design, one on trade-offs and
        critique, and a lab for building and deploying.
\end{itemize}

You should be comfortable reading Python and have used a large language model
through an API. You do not need prior experience with agent frameworks, MCP, or
cloud deployment; the book introduces them as the cases require them.

\section*{How to read this book}

The cases can be read in order, and I recommend doing so the first time
through, because the early cases introduce design patterns that later cases
combine. After that, each case stands on its own, and readers interested in a
particular industry can go straight to it.

Above all, I encourage you to draw before you read. At the start of each case,
stop after the business problem, take a blank sheet of paper, and sketch your
own design. Then read on and compare yours with mine. Where we differ, ask
whether your design is worse, or simply makes a different trade-off. Often it
will be the latter, and learning to tell the two apart is the point of the
exercise.

\section*{The AI companion}

Every reader of this book receives free access to an AI companion: an assistant
built to work alongside these chapters. You can use it to talk through a case,
to critique a design you have drawn, to ask ``what if'' questions about
alternative architectures, or to get unstuck in a lab. Think of it as a study
partner who has read the whole book and never tires of being asked why.

The companion is also a small demonstration of the book's subject. It is itself
a designed agentic solution, and it is worth asking, as you use it, how you
would have designed it differently.

\section*{Companion materials}

All code for the cases, along with the diagrams and deployment
configurations, is available in the book's companion repository. Because the
tools in this field change quickly, the repository will be kept up to date even
when the printed text cannot be.

% TODO: insert repository URL and AI companion URL.

\section*{Acknowledgments}

This book grew out of questions from my students, and my first thanks go to
them. Their confusion about where to begin was the most useful feedback any
author could ask for.

% TODO: add personal acknowledgments (reviewers, colleagues, publisher, family).

I hope this book helps you move from knowing the parts to designing the whole,
and from the whiteboard to production.

\vspace{2em}
\begin{flushright}
Chris Kuo\\
% TODO: place and date
\end{flushright}
